\section{A sharp stability scale at the circular configuration}
\label{sec:v4-coalescence}
The exact infinite-sequence identification in
Theorem~\ref{thm:g-inverse} does not resolve conditioning at a repeated
node. We now give both a sequence-level estimate for the coalescing
path and a matching positive estimate relative to the circular
configuration. The latter is explicitly a reference-point estimate;
it is not asserted for every pair of noncircular triples.

\subsection{An observable dispersion statistic}
Use the equal-gap family \eqref{eq:g-isogap-family}, with $R$ fixed
and $g=1-2R$. Define, from its first three positive amplitudes,
\begin{equation}\label{eq:v4-dispersion-data}
 m=\frac{2C_2}{C_1},\qquad
 J=\frac{C_3}{C_1}-f(m),\qquad f(x)=\frac{x^2}{4-x^2}.
\end{equation}
For the circular member let
$\kappa_*=1/R$ and $x_*=(1+g\kappa_*)^{-1}$.

\begin{theorem}[Three-amplitude circular-reference stability]
\label{thm:v4-radial-stability}
On a sufficiently small compact positive-curvature neighborhood in the
fixed-$R$ family, $J\ge0$ and there are constants $c,C>0$ such that
\begin{equation}\label{eq:v4-radial-comparison}
 c\left(|m-x_*|+\sqrt J\right)
 \le \max_{0\le r\le2}|\kappa_r-\kappa_*|
 \le C\left(|m-x_*|+\sqrt J\right).
\end{equation}
In particular, if $C_j^*$ are the circular amplitudes and
$\delta=\max_{1\le j\le3}|C_j-C_j^*|\le1$, then
\begin{equation}\label{eq:v4-radial-holder}
 \operatorname{dist}_{\rm match}(\kappa,\kappa_*)
                                                    \le C\sqrt\delta.
\end{equation}
The statistic $J$ vanishes exactly when the three contact curvatures
coincide. It measures their dispersion without a supplied area
normalizer or channel labels.
\end{theorem}
\begin{proof}
Set $x_r=\operatorname{sech}\gamma_r=(1+g\kappa_r)^{-1}$ and
\[
 p_r=\frac{\csch\gamma_r}{\sum_s\csch\gamma_s}.
\]
The $x_r$ lie in a fixed compact interval inside $(0,1)$, while
$p_r\ge p_*>0$ and $\sum p_r=1$. The identities
$\sinh(2\gamma)=2\sinh\gamma\cosh\gamma$ and
$\sinh(3\gamma)=\sinh\gamma(4\cosh^2\gamma-1)$ give
\begin{equation}\label{eq:v4-jensen}
 m=\sum_rp_rx_r,\qquad J=\sum_rp_rf(x_r)-f(m).
\end{equation}
On the common interval,
\[
 f''(x)=\frac{8(4+3x^2)}{(4-x^2)^3}
\]
has positive lower and finite upper bounds. Taylor's integral formula
at $m$, followed by $\sum p_r(x_r-m)=0$, therefore gives
\begin{equation}\label{eq:v4-jensen-bounds}
 \frac{c_f}{2}\sum_rp_r(x_r-m)^2\le J
       \le\frac{C_f}{2}\sum_rp_r(x_r-m)^2.
\end{equation}
This proves nonnegativity and the equality criterion. Since all
weights are bounded below, the square root in
\eqref{eq:v4-jensen-bounds} is comparable to
$\max_r|x_r-m|$. Also
\[
 |m-x_*|\le\max_r|x_r-x_*|,\qquad
 \max_r|x_r-x_*|\le |m-x_*|+\max_r|x_r-m|.
\]
The reverse comparison follows from
$\max_r|x_r-m|\le2\max_r|x_r-x_*|$.
Finally the scalar change $\kappa\mapsto(1+g\kappa)^{-1}$ is
bi-Lipschitz on the compact interval. This proves
\eqref{eq:v4-radial-comparison}.

All amplitude denominators $C_1$ are bounded away from zero. At the
circle $m=x_*$ and $J=0$. The two ratios in
\eqref{eq:v4-dispersion-data} are Lipschitz in the first three
amplitudes on this set, and $f$ has bounded derivative. Hence
$|m-x_*|\le C\delta$ and $0\le J\le C\delta$.
Substitution gives \eqref{eq:v4-radial-holder}. Because the reference
triple has all entries equal, its matching distance is exactly the
maximum coordinate distance, independent of ordering.
\end{proof}

The theorem quantifies distance to a specified circular table. It
does not recover an arbitrary unordered triple from only three
numbers, nor does it claim a one-half-H\"older inverse between every
two nearby noncircular triples. The exact infinite-sequence theorem
and its multiplicity recovery remain separate statements.

\subsection{The infinite absolute-amplitude topology}
For a fixed $a\ge0$ write
\begin{equation}\label{eq:v4-sequence-norm}
 \|D\|_a=\sup_{j\ge1}e^{aj}|D_j|.
\end{equation}
Restrict the coalescing path of
Proposition~\ref{prop:v3-area-coalescence} to a small closed interval
on which all exponents satisfy $\gamma_r(s)\ge\gamma_->a$.

\begin{theorem}[Sharp circular-reference exponent for the full sequence]
\label{thm:v4-sequence-stability}
Along $\alpha=\zeta=0$, $\beta=s$ one has
\begin{equation}\label{eq:v4-sequence-path}
 \|C(s)-C(0)\|_a\le B_a s^2,
 \qquad
 \operatorname{dist}_{\rm match}(\kappa(s),\kappa(0))\asymp|s|.
\end{equation}
For every table in a sufficiently small fixed-$R$ neighborhood of the
circle, the positive estimate
\begin{equation}\label{eq:v4-sequence-upper}
 \operatorname{dist}_{\rm match}(\kappa,\kappa_*)
                   \le C_a\|C-C^*\|_a^{1/2}
\end{equation}
holds when the right side is small. No exponent larger than one half
can replace it throughout that neighborhood, even when the entire
sequence is observed. This sharpness concerns the absolute-error
norm \eqref{eq:v4-sequence-norm}, with a strict exponential margin.
\end{theorem}
\begin{proof}
On the indicated path the three radii are $R+36s,R-18s,R-18s$,
and $A(s)=A(0)+45\pi s^2/4$. With
$\Phi_j(r)=\csch(j\arcosh(1+g/r))$ the amplitude is
\[
 C_j(s)=\frac{\Phi_j(R+36s)+2\Phi_j(R-18s)}{A(s)}.
\]
Both numerator and denominator have zero first derivative at zero.
On the compact interval the derivatives through order two of the
radius-to-exponent map and of $A^{-1}$ are bounded. Differentiate
\[
 \csch(j\gamma)=\frac{2e^{-j\gamma}}{1-e^{-2j\gamma}}
\]
twice. The denominator is bounded below by $1-e^{-2\gamma_-}$;
each differentiation contributes at most a factor $C(1+j)$.
Consequently
\begin{equation}\label{eq:v4-sequence-second}
 \sup_{|s|\le s_*}|C_j''(s)|
                         \le B(1+j)^2e^{-j\gamma_-}.
\end{equation}
Taylor's integral remainder and $C_j'(0)=0$ imply the first part of
\eqref{eq:v4-sequence-path}, since
$\sup_j(1+j)^2e^{-(\gamma_--a)j}<\infty$.
The curvature displacement is comparable to $|s|$ because reciprocal
radius has a derivative bounded above and away from zero there.

For a general nearby table, the first three absolute amplitude errors
are bounded by \eqref{eq:v4-sequence-norm}. Thus
Theorem~\ref{thm:v4-radial-stability} gives
\eqref{eq:v4-sequence-upper}; no unstable truncation of a M\"obius
inversion is used. If an exponent $\theta>1/2$ were valid in its place,
the specified path would give
$b|s|\le L B_a^\theta|s|^{2\theta}$ for arbitrarily small nonzero
$s$, which is impossible. This proves sharpness. In particular,
\eqref{eq:v4-sequence-upper} also supplies a quadratic lower bound
for the sequence discrepancy along that path; the order in
\eqref{eq:v4-sequence-path} is not merely a one-sided diagnosis.
\end{proof}

Exact identification and stability in an absolute-amplitude norm are
different questions. The two theorems give a positive sharp scale for
one specified singular reference, while the channel-resolved endpoint
observation remains Lipschitz there because it retains different data.
None of these assertions treats relative error at arbitrarily small
amplitudes as equivalent to \eqref{eq:v4-sequence-norm}.
